\documentclass[aspectratio=169,11pt]{beamer}

\usetheme{Madrid}
\usecolortheme{default}
\usefonttheme{professionalfonts}
\setbeamertemplate{navigation symbols}{}
\setbeamertemplate{footline}[frame number]

\usepackage{amsmath, amssymb, booktabs, array}

\title{Behavioral Public Policy in Labor Markets}
\subtitle{Behavioral Labor -- Week 9}
\author{Teaching deck draft}
\date{}

\begin{document}

\begin{frame}
  \titlepage
\end{frame}

\begin{frame}{Central question}
\begin{itemize}
    \item How do behavioral frictions change the design, implementation, targeting, equilibrium effects, and welfare analysis of labor-market policies?
    \item The policy object is not only generosity or tax rates.
    \item Workers respond to perceived incentives, program complexity, defaults, trust, deadlines, and learning environments.
    \item The labor object stays central: earnings, claiming, search, training, saving, and work incentives.
\end{itemize}
\end{frame}

\begin{frame}{Bridge from Weeks 1--7}
\small
\begin{tabular}{p{0.24\linewidth} p{0.60\linewidth}}
\toprule
Earlier weeks & Week 9 policy task \\
\midrule
Preferences & design around present bias, commitment, inertia, and follow-through \\
Beliefs & correct mistaken returns, eligibility, and schedule perceptions \\
Attention & make incentives and rules salient at the relevant decision moment \\
Complexity & treat forms, interfaces, and burden as economic margins \\
Firms and norms & study delivery through employers, counselors, and local environments \\
Methods & match policy channels to identification and welfare tools \\
\bottomrule
\end{tabular}

\vspace{0.7em}
Week 9 turns the course toolkit into a behavioral labor-policy framework.
\end{frame}

\begin{frame}{Six roles of behavioral frictions}
\small
\begin{tabular}{p{0.29\linewidth} p{0.56\linewidth}}
\toprule
Role & Policy question \\
\midrule
Bias to correct & what misunderstanding, inattention, or self-control wedge blocks welfare-improving action? \\
Lever to harness & can defaults, reminders, commitment, or framing improve outcomes? \\
Implementation friction & is take-up limited by hassle, stigma, trust, or procedural complexity? \\
Dynamic learning margin & do workers learn schedules, rules, or procedures over time? \\
Equilibrium modifier & how do employers, preparers, counselors, and peers mediate response? \\
Welfare complication & which choices are distorted, and whose benchmark matters? \\
\bottomrule
\end{tabular}
\end{frame}

\begin{frame}{Standard policy versus behavioral policy}
\[
p = (\tau, n, s, d)
\]

\small
\begin{tabular}{p{0.18\linewidth} p{0.63\linewidth}}
\toprule
Component & Interpretation \\
\midrule
$\tau$ & taxes, subsidies, benefit rates, budget-set parameters \\
$n$ & information, salience, notices, rule communication \\
$s$ & simplification, administrative burden, claiming support \\
$d$ & defaults, commitment, active-choice architecture \\
\bottomrule
\end{tabular}

\vspace{0.7em}
Standard policy often emphasizes $\tau$. Behavioral labor policy asks which component of $p$ workers actually perceive and can act on.
\end{frame}

\begin{frame}{Worker choice under behavioral policy}
\[
a_i(p) \in \arg\max_{a \in \mathcal{A}} \tilde U_i(a; p, b_i) - K_i(a; p)
\]

\begin{itemize}
    \item $a_i(p)$: hours, earnings, search, training, claiming, saving, or work after disability.
    \item $b_i$: beliefs, attention, sophistication, perceived eligibility, trust, or local knowledge.
    \item $K_i(a;p)$: procedural, cognitive, time, stigma, or coordination costs.
    \item Policy can move objective incentives, perceived incentives, action costs, defaults, or learning.
\end{itemize}
\end{frame}

\begin{frame}{Biases as targets of correction}
\begin{itemize}
    \item EITC response may be weak because workers do not know the phase-in, plateau, or phaseout.
    \item Social Security or disability rules may be misperceived, so claiming and work incentives are misunderstood.
    \item Training after displacement may be delayed because costs are immediate and returns are uncertain.
    \item Corrective design asks: what would workers choose with better information, lower hassle, or less distortion?
    \item Empirical anchors: EITC knowledge, tax-benefit learning, Social Security information.
\end{itemize}
\end{frame}

\begin{frame}{Biases as tools and levers}
\begin{itemize}
    \item Defaults can turn inertia into retirement saving.
    \item Active choice can force attention without choosing for the worker.
    \item Reminders can shift deadline attention and follow-through.
    \item Commitment devices can help when short-run costs block long-run investments.
    \item Large behavior change is not automatically large welfare gain.
\end{itemize}
\end{frame}

\begin{frame}{Implementation and take-up}
\[
D_i(p) = \mathbf{1}\!\left\{\tilde V_i(p; b_i) - C_i(p) \ge 0\right\}
\]

\begin{itemize}
    \item Take-up is a labor-policy margin, not a clerical detail.
    \item Low participation can reflect low awareness, misunderstood eligibility, stigma, distrust, or paperwork.
    \item Policy can change $D_i(p)$ through salience, simplification, reminders, assistance, and timing.
    \item The main outcome may be claiming or completion, not hours or earnings.
\end{itemize}
\end{frame}

\begin{frame}{Administrative burden as economics}
\small
\begin{tabular}{p{0.25\linewidth} p{0.28\linewidth} p{0.30\linewidth}}
\toprule
Policy setting & Behavioral margin & Outcome object \\
\midrule
Tax credits & eligibility and schedule knowledge & filing, claiming, earnings response \\
UI & search rules and deadlines & applications, hazards, reemployment \\
Training & enrollment and completion burden & take-up, credits, credentials \\
Retirement & defaults and payroll menus & participation, contribution rates \\
Disability & earnings rules and perceived risk & claiming, work, exits \\
\bottomrule
\end{tabular}
\end{frame}

\begin{frame}{Dynamic learning and persistence}
\[
b_{i,t+1} = B\!\left(b_{it},\, m_{it},\, x_{i,t+1},\, p_t\right)
\]

\begin{itemize}
    \item Workers learn schedules, rules, claiming procedures, and returns over time.
    \item Information acquisition can be endogenous to stakes, deadlines, peers, and intermediaries.
    \item Short-run effects may understate long-run learning or overstate temporary salience.
    \item Dynamic designs need panels, event studies, distributed lags, or repeated exposure measures.
\end{itemize}
\end{frame}

\begin{frame}{Policy-effect channels}
\[
\frac{d a_i}{d p}
=
\frac{\partial a_i}{\partial \tau}\frac{d\tau}{dp}
+
\frac{\partial a_i}{\partial n}\frac{dn}{dp}
+
\frac{\partial a_i}{\partial s}\frac{ds}{dp}
+
\frac{\partial a_i}{\partial d}\frac{dd}{dp}
\]

\begin{itemize}
    \item The decomposition is a teaching device, not a literal universal derivative.
    \item It keeps incentives, salience, simplification, and defaults on the same page.
    \item Good empirical design tries to isolate or diagnose the active channel.
\end{itemize}
\end{frame}

\begin{frame}{Firms, intermediaries, and equilibrium}
\begin{itemize}
    \item Employers set payroll menus, retirement defaults, training supports, and communication.
    \item Tax preparers translate tax-credit schedules and deadlines.
    \item Counselors and caseworkers shape UI, training, disability, and claiming behavior.
    \item Peers and neighborhoods diffuse knowledge, norms, and partial information.
    \item Equilibrium response can change incidence, pass-through, and scalability.
\end{itemize}
\end{frame}

\begin{frame}{Welfare, targeting, and ambiguity}
\[
W(p) = \int \Psi_i\!\left(a_i(p), a_i^\star(p), p\right)\, dF(i)
\]

\begin{itemize}
    \item Revealed preference is harder to use when choices may be distorted.
    \item $a_i^\star(p)$ could mean full-information choice, low-hassle choice, sophisticated choice, or model benchmark.
    \item Paternalism is a risk if real preferences are misclassified as mistakes.
    \item Targeting can prioritize high mistakes, high gains, low capacity, or equity concerns.
\end{itemize}
\end{frame}

\begin{frame}{Methods bridge}
\scriptsize
\begin{tabular}{p{0.31\linewidth} p{0.30\linewidth} p{0.24\linewidth}}
\toprule
Policy channel & Common methods & Key diagnostic \\
\midrule
Information, reminders & field experiments, ANCOVA & baseline knowledge and heterogeneity \\
Take-up, claiming & extensive-margin effects, hazards & timing and application completion \\
Learning & panels, event studies, lags & persistence versus temporary salience \\
Defaults, menus & treatment comparisons, bunching, calibration & passive choice and welfare benchmark \\
Peer or local knowledge & local interactions, spillover designs & diffusion and interference \\
Welfare, targeting & sufficient statistics, structural models & benchmark action and distribution \\
\bottomrule
\end{tabular}
\end{frame}

\begin{frame}{Frontier directions}
\begin{itemize}
    \item Separate information, salience, procedural knowledge, trust, and confidence.
    \item Study persistence, decay, repeated exposure, and endogenous information acquisition.
    \item Treat digital administration as a behavioral environment, not neutral plumbing.
    \item Build intermediaries and firm response into policy design.
    \item Make welfare criteria explicit when choices may be distorted.
\end{itemize}
\end{frame}

\begin{frame}{Week 9 lab}
\begin{itemize}
    \item Reproduce: synthetic take-up and reminder design anchored on Bhargava--Manoli.
    \item Diagnose: targeted friction, correction versus harnessing, main margin, welfare benchmark, intermediary response.
    \item Transfer: EITC knowledge and local labor-supply response anchored on Chetty--Friedman--Saez.
    \item Challenge: defaults and welfare, or dynamic tax-benefit learning.
\end{itemize}
\end{frame}

\begin{frame}{Bridge from Week 8}
\begin{itemize}
    \item Week 8 made firms, platforms, markets, and institutions strategic actors.
    \item Week 9 asks how policy should be designed once those responses are present.
    \item Implementation tools may change incidence when employers and intermediaries adapt.
    \item Welfare requires both distorted-choice benchmarks and equilibrium incidence.
\end{itemize}
\end{frame}

\begin{frame}{Takeaways}
\begin{itemize}
    \item Behavioral labor policy is broader than nudges.
    \item Incentives, information, simplification, defaults, and learning are joint design objects.
    \item Implementation and take-up are economic margins.
    \item Intermediaries can amplify, mute, or redirect policy effects.
    \item Welfare analysis requires an explicit benchmark for distorted choice.
\end{itemize}
\end{frame}

\end{document}
